% ============================================================
%  MG-HGNN: Modality-Gated Heterogeneous Graph Neural Networks
%  for Student Academic Performance Prediction
%  Conference Paper — ACM SIGCONF Format
% ============================================================

\documentclass[sigconf]{acmart}

% ---- Packages ----
\usepackage{amsmath}
\usepackage{amssymb}
\usepackage{graphicx}
\usepackage{booktabs}
\usepackage{hyperref}
\usepackage{algorithm2e}
\usepackage{tikz}
\usepackage{lipsum}

% ---- ACM Metadata (update before submission) ----
\acmConference[KDD '26]{ACM SIGKDD Conference on Knowledge Discovery and Data Mining}
              {August 2026}{Barcelona, Spain}
\acmDOI{}
\acmISBN{}

% ============================================================
\begin{document}

% ---- Title ----
\title{MG-HGNN: Modality-Gated Heterogeneous Graph Neural Networks\\
       for Student Academic Performance Prediction}

% ---- Authors (update affiliations) ----
\author{[Author Name]}
\affiliation{%
  \institution{[Your University / Institution]}
  \city{[City]}
  \country{[Country]}
}
\email{[email@university.edu]}

% ---- CCS Concepts and Keywords ----
\begin{CCSXML}
<ccs2012>
  <concept>
    <concept_id>10010147.10010257.10010293.10010294</concept_id>
    <concept_desc>Computing methodologies~Neural networks</concept_desc>
    <concept_significance>500</concept_significance>
  </concept>
  <concept>
    <concept_id>10010147.10010178</concept_id>
    <concept_desc>Computing methodologies~Artificial intelligence</concept_desc>
    <concept_significance>300</concept_significance>
  </concept>
</ccs2012>
\end{CCSXML}

\ccsdesc[500]{Computing methodologies~Neural networks}
\ccsdesc[300]{Computing methodologies~Artificial intelligence}

\keywords{heterogeneous graph neural networks, student performance prediction,
          multi-modal learning, modality gating, educational data mining,
          learning analytics}

% ============================================================
%  ABSTRACT
% ============================================================
\begin{abstract}
Predicting student academic performance is a critical task in educational data
mining, enabling early identification of at-risk learners and timely pedagogical
intervention. Recent advances in graph neural networks (GNNs) have demonstrated
strong potential by explicitly modelling relational structures among students,
courses, and learning activities. However, existing heterogeneous GNN (HGNN)
approaches share a fundamental limitation: they collapse all available modalities
--- structured academic records, textual content, and behavioral interaction logs
--- into a single, unified node embedding. This flattening ignores the fact that
different edge relation types inherently carry different modal signals; for
instance, a student's behavioral clickstream is more informative when propagating
messages through a \emph{resource-access} edge than through a
\emph{peer-collaboration} edge, where textual communication patterns are
dominant.

In this paper, we propose \textbf{MG-HGNN}, a \textit{Modality-Gated
Heterogeneous Graph Neural Network} that addresses this gap. Each node in the
graph maintains three \emph{separate} modality embeddings --- a structured
embedding produced by an MLP encoder, a textual embedding from a fine-tuned BERT
encoder, and a behavioral embedding from a bidirectional LSTM encoder. The core
innovation is a lightweight, learnable \emph{modality gate}: for each edge
relation type $r$, a dedicated softmax gate $\boldsymbol{\alpha}^{(r)} =
[\alpha_s^{(r)},\, \alpha_t^{(r)},\, \alpha_b^{(r)}]$ is learned end-to-end,
dynamically controlling how much each modality contributes to message passing
along that specific relation. The gate mechanism is conditioned exclusively on
the relation type, imposing no additional overhead at the node level.

We evaluate MG-HGNN on two publicly available educational datasets --- MOOC-Cube
and the Open University Learning Analytics Dataset (OULAD) --- across three
complementary prediction tasks: final grade regression, dropout risk
classification, and engagement-level prediction. Experimental results show that
MG-HGNN consistently outperforms strong baselines including HAN, HGT, R-GCN,
and multi-topology GNN variants. Ablation studies confirm that edge-conditioned
gating is the key factor driving the performance gain, and post-training
analysis of learned gate weights validates our core hypothesis: the model
autonomously discovers that behavioral signals are most relevant for
resource-access edges, textual signals for collaboration edges, and structured
signals for enrollment edges --- without any manual specification of these
inductive biases.
\end{abstract}

\maketitle

% ============================================================
%  1. INTRODUCTION
% ============================================================
\section{Introduction}
\label{sec:intro}
\lipsum[1-2]
% TODO: Expand with motivation, problem statement, and contributions.

% ============================================================
%  2. RELATED WORK
% ============================================================
\section{Related Work}
\label{sec:related}

The prediction of student academic performance has evolved substantially over the
past decade, progressing from classical machine learning models operating on
independent per-student feature vectors toward relational deep learning methods
that explicitly encode the complex network structure of educational ecosystems.
This section surveys the most relevant families of approaches, situating our
proposed MG-HGNN within the existing landscape.

\subsection{Graph Neural Networks for Student Performance Prediction}

Early graph-based approaches demonstrated that modelling student similarity as an
explicit graph structure yields systematic gains over flat-feature baselines.
\citet{li2022studygnn} proposed Study-GNN, a multi-topology GNN (MTGNN) in which
multiple student similarity graphs are constructed using distinct distance
metrics. An attention mechanism then fuses these topologies, casting performance
prediction as semi-supervised node classification where each node represents a
student and labels encode pass, fail, or withdrawal outcomes. Study-GNN was shown
to outperform support vector machines, logistic regression, shallow neural
networks, and vanilla graph convolutional networks, validating the utility of
multi-view relational modelling.

Building on the insight that student--student similarity alone is insufficient,
\citet{huang2024dual} introduced APP-DGNN, a dual graph neural network that
constructs two complementary graphs: one derived from interaction logs capturing
engagement dynamics over time, and another from student attributes such as
demographics and academic background. Local representations from the interaction
graph and global representations from the attribute graph are fused for
downstream prediction, achieving approximately 84\% accuracy on pass/fail
classification and approximately 90\% on pass/withdraw tasks on a public MOOC
dataset. This dual-graph paradigm highlights a key architectural principle:
different information sources warrant distinct graph topologies rather than a
single aggregated representation.

\subsection{Temporal and Dynamic Graph Models}

A recurring limitation of static graph methods is their inability to capture the
evolution of student behavior and performance across time. \citet{huang2024temporal}
addressed this with APP-TGN, a temporal graph network (TGN) in which student
activity logs are encoded as a dynamic graph where edges carry timestamps. The
model combines low-pass and high-pass spectral filters with global neighborhood
sampling and multi-head attention, significantly improving both overall predictive
accuracy and early at-risk detection in massive open online course (MOOC)
settings. Similarly, \citet{zhang2025gnntinet} demonstrated the value of
incorporating temporal context within a richer hybrid architecture, embedding GNN
layers inside a GNN-Transformer-InceptionNet pipeline for multi-label student
performance prediction, reporting up to 98.5\% accuracy on a single-institution
dataset.

Despite these advances, several dual and multi-graph designs still do not fully
exploit long-range temporal dependencies \citep{huang2024dual, huang2024temporal,
balachandar2024mspp}, motivating continued work on temporally expressive
architectures. Our proposed MG-HGNN addresses a complementary and previously
unexplored axis: rather than improving temporal modelling, we improve
\emph{modal selectivity} during heterogeneous message passing.

\subsection{Hypergraph and Higher-Order Models}

Beyond pairwise edges, \citet{yang2024hypergraph} proposed FD-HGNN, a framelet-
based dual hypergraph neural network for student performance prediction.
Framelet transforms are used to produce low-pass and high-pass hypergraphs over
students, and a dual-channel attention mechanism models higher-order, multi-way
relationships among student groups rather than individual pairs. Evaluated across
four datasets, FD-HGNN outperforms both traditional machine learning methods and
standard GNNs. This line of work underscores the importance of capturing
structural complexity beyond simple pairwise student--course relations; however,
hypergraph constructions do not directly address the problem of multi-modal
feature heterogeneity within each node.

\subsection{Multi-source and Heterogeneous Data Fusion}

Several works have emphasised the need to handle sparse, multi-source data
including behavioral logs, demographic features, emotional indicators, and prior
academic records more effectively than flat concatenation allows
\citep{huang2024dual, zhang2025gnntinet, huang2024temporal, li2022studygnn,
yang2024hypergraph, balachandar2024mspp}. \citet{balachandar2024mspp} proposed
the Multi-Dimensional Student Performance Prediction model (MSPP), which
incorporates advanced graph neural network layers on heterogeneous academic data,
achieving approximately 76\% accuracy on complex multi-class classification tasks.
This work demonstrates that even with heterogeneous input, current methods still
rely on pre-fusion of modalities before graph operations, limiting the
expressiveness of relational message passing.

Collectively, these contributions establish graph-based deep learning as the
dominant paradigm for student performance modelling, with reported accuracy
ranging from approximately 76\% on complex multi-class tasks
\citep{balachandar2024mspp} to 84--90\% on binary outcomes \citep{huang2024dual}
and up to 98--99\% in richer, typically single-institution settings
\citep{zhang2025gnntinet}. Shared limitations across this body of work include
high computational cost for complex GNN operators on large student cohorts, and
questions about cross-institutional generalizability
\citep{huang2024dual, li2022studygnn, zhang2025gnntinet, huang2024temporal,
yang2024hypergraph, balachandar2024mspp}.

\subsection{Gap Addressed by This Work}

Despite the breadth of existing approaches, a fundamental assumption has remained
unexamined: that all relation types in the educational graph should draw on the
same mixture of input modalities when aggregating neighborhood information. No
existing HGNN framework conditions its feature fusion strategy on the type of
edge being traversed. This is the precise gap that MG-HGNN fills. By maintaining
separate per-modality embeddings at each node and learning a distinct, softmax-
normalized modality gate $\boldsymbol{\alpha}^{(r)}$ for each edge relation type
$r \in \mathcal{R}$, MG-HGNN allows the graph itself to discover which
information sources are most relevant for each relational context --- without
requiring manual specification of these inductive biases.

% ============================================================
%  3. PROBLEM DEFINITION
% ============================================================
\section{Problem Definition}
\label{sec:problem}
\lipsum[1]
% TODO: Formal graph definition G = (V, E, T_v, T_e), node/edge types, tasks.

% ============================================================
%  4. METHODOLOGY
% ============================================================
\section{Methodology: MG-HGNN}
\label{sec:method}
\lipsum[1-2]
% TODO: Modality encoders, gate formulation, HGT backbone, prediction heads.

% ============================================================
%  5. EXPERIMENTS
% ============================================================
\section{Experiments}
\label{sec:experiments}
\lipsum[1]
% TODO: Datasets (MOOC-Cube, OULAD), baselines, RQ1-RQ4, results tables.

% ============================================================
%  6. CONCLUSION
% ============================================================
\section{Conclusion}
\label{sec:conclusion}
\lipsum[1]
% TODO: Summary of contributions, limitations, and future work.

% ============================================================
%  REFERENCES (BibTeX)
% ============================================================
\bibliographystyle{ACM-Reference-Format}

\begin{thebibliography}{9}

\bibitem{balachandar2024mspp}
Balachandar, V., \& Venkatesh, K. (2024).
\newblock A multi-dimensional student performance prediction model ({MSPP}):
  An advanced framework for accurate academic classification and analysis.
\newblock \emph{MethodsX}, 14.
\newblock \url{https://doi.org/10.1016/j.mex.2024.103148}

\bibitem{huang2024dual}
Huang, Q., \& Zeng, Y. (2024).
\newblock Improving academic performance predictions with dual graph neural
  networks.
\newblock \emph{Complex \& Intelligent Systems}, 10, 3557--3575.
\newblock \url{https://doi.org/10.1007/s40747-024-01344-z}

\bibitem{huang2024temporal}
Huang, Q., \& Chen, J. (2024).
\newblock Enhancing academic performance prediction with temporal graph
  networks for massive open online courses.
\newblock \emph{Journal of Big Data}, 11, 1--26.
\newblock \url{https://doi.org/10.1186/s40537-024-00918-5}

\bibitem{li2022studygnn}
Li, M., Wang, X., Wang, Y., Chen, Y., \& Chen, Y. (2022).
\newblock {Study-GNN}: A novel pipeline for student performance prediction
  based on multi-topology graph neural networks.
\newblock \emph{Sustainability}.
\newblock \url{https://doi.org/10.3390/su14137965}

\bibitem{li2022campus}
Li, X., Zhang, Y., Cheng, H., Li, M., \& Yin, B. (2022).
\newblock Student achievement prediction using deep neural network from
  multi-source campus data.
\newblock \emph{Complex \& Intelligent Systems}, 8, 5143--5156.
\newblock \url{https://doi.org/10.1007/s40747-022-00731-8}

\bibitem{yang2024hypergraph}
Yang, Y., Shi, J., Li, M., \& Fujita, H. (2024).
\newblock Framelet-based dual hypergraph neural networks for student
  performance prediction.
\newblock \emph{International Journal of Machine Learning and Cybernetics},
  15, 3863--3877.
\newblock \url{https://doi.org/10.1007/s13042-024-02124-4}

\bibitem{zhang2025gnntinet}
Zhang, X., Zhang, Y., Chen, A., Yu, M., \& Zhang, L. (2025).
\newblock Optimizing multi-label student performance prediction with
  {GNN-TINet}: A contextual multidimensional deep learning framework.
\newblock \emph{PLOS ONE}, 20.
\newblock \url{https://doi.org/10.1371/journal.pone.0314823}

\end{thebibliography}

\end{document}
